A diferença entre os valores de \gls{rmse} simulado e experimental é atribuída a
fatores não contemplados pelo modelo \gls{fopdt}: ruído de medição, atrasos de
transporte e dinâmica térmica não modelada, além da saturação de potência do
aquecedor. Contribui ainda a estrutura do problema, em que o controlador é
realimentado pelo sensor~1 enquanto o erro é avaliado no sensor~2, cuja
diferença de localização introduz dinâmica adicional não capturada pelo modelo.
Apesar dessas divergências, a proximidade entre os valores de \gls{rmse} indica
que o modelo representa adequadamente a planta e que a estratégia de controle
mantém bom desempenho experimental.